Mit der spezifischen Wärmekapazität von Wasser $c = \cWO$ (FoTaBe S.~191) und $\Delta T = \dTO$ ist die freigesetzte Wärme
\begin{align}
 Q &= \QmF = \cW\cdot\mM\cdot\dT = \Qm \approx \result{\QmS{2}}
\end{align}
Der Weltenergiebedarf beträgt $E_\text{W} = \EwO = \Ew$ pro Jahr. Die Wärme würde also für
\begin{align}
 \frac{Q}{E_\text{W}} &= \nJF = \frac{\Qm}{\Ew} = \nJ \approx \result{\nJP{0}{2}}
\end{align}
Jahre reichen.

Wärme lässt sich aber nur mit einer Wärmekraftmaschine in Arbeit umwandeln, und diese braucht ein kälteres Reservoir, an das sie Abwärme abgeben kann. Ein solches ist nicht vorhanden: Das Meer ist selbst das kälteste grosse Reservoir. Selbst wenn eines zur Verfügung stünde, wäre bei einer Temperaturdifferenz von nur $\Delta T = \dTO$ bei $T \approx \ToO$ der Wirkungsgrad höchstens
\begin{align}
 \eta_\text{C} &= \etCF = \frac{\dT}{\To} = \etC
\end{align}
Man müsste also riesige Wassermengen umwälzen, um wenig Arbeit zu gewinnen. Nach dem zweiten Hauptsatz kann die Wärme eines Reservoirs ohne Temperaturgefälle nicht in Arbeit umgewandelt werden.
